% Part: proof-theory
% Chapter: natural-deduction
% Section: quantifiers

\documentclass[../../../include/open-logic-section]{subfiles}

\begin{document}

\olfileid{pt}{nat}{qua}

\olsection{नियमित \usetoken{P}{proof} आणि आदेशन}

संख्यापकांच्या नियमांमध्ये
~\Elim{\lexists} आणि
~\Intro{\lforall} यांना
लागणाऱ्या “आयगेन चराच्या
अटींमुळे” काही गुंतागुंत येते.
या अनुमानांतील
!!{constant}~$c$ पुन्हा
वापरला जाणार नाही याची
खात्री केल्यास बहुतांश
गुंतागुंत हाताळता येते.
पुढील व्याख्या करू.

\begin{defn}
नैसर्गिक निगमनातील
!!{proof}~$\delta$
\emph{नियमित} असते, जर
\begin{enumerate}
  \item कोणतेही आयगेन
  चर~$a$ एकाहून अधिक
  \Elim{\lexists} किंवा
  ~\Intro{\lforall} अनुमानांचे
  आयगेन चर म्हणून वापरलेले
  नसेल; आणि
  \item $a$ हे
  \Elim{\lexists} किंवा
  ~\Intro{\lforall} अनुमानाचे
  आयगेन चर असेल, तर ते
  अनुक्रमे गौण आधारविधानाकडे
  किंवा आधारविधानाकडे
  जाणाऱ्या निकटच्या
  उप-!!{proof} मध्येच
  येत असेल.
\end{enumerate}
\end{defn}

\begin{lem}\ollabel{lem:subst-regular}
$\delta$ ही $!C$ मध्ये
संपणारी नियमित सिद्धता असू द्या.
$c$ हा~$\delta$ मध्ये
आयगेन चर म्हणून न वापरलेला
!!a{constant} असू द्या;
आणि~$t$ हे $\delta$
चे कोणतेही आयगेन चर नसलेले
पद असू द्या. मग~$\delta$
मधील $c$ च्या प्रत्येक
उपस्थितीच्या जागी $t$
बसवून मिळणारी
$\Subst{\delta}{t}{c}$
ही नियमित !!{proof} आहे.
$\Subst{\delta}{t}{c}$
चे अंतिम-!!{formula}
$\Subst{!C}{t}{c}$ असते.
$!A^x$ हे~$\delta$
चे उघडे गृहीतक असेल,
तर $\Subst{!A}{t}{c}^x$
हे $\Subst{\delta}{t}{c}$
चे उघडे गृहीतक असते.
\end{lem}

\begin{proof}
  $\pheight{\delta}$ वर
  विगमन करू.
  $\pheight{\delta} = 0$
  असेल, तर सिद्धतेत केवळ
  गृहीतक~$!A^x$ असते.
  म्हणून $\Subst{\delta}{t}{c}$
  हीच~$\Subst{!A}{t}{c}^x$
  असते.

  $\pheight{\delta} > 0$
  असेल, तर~$\delta$
  मधील अंतिम अनुमानानुसार
  प्रसंग पाडू.
  विधानीय संयोजकांच्या
  नियमांचे प्रसंग सोपे
  असून सरावासाठी ठेवले आहेत.
  $\delta$ चा शेवट
  संख्यापक अनुमानात होतो
  ते प्रसंग पाहू.
  \begin{enumerate}
    \item अंतिम अनुमान
    $\Intro{\lforall}$
    असेल, तर~$\delta$
    चे रूप असे आहे:
    \[
    \AxiomC{}
    \RightLabel{$\delta'$}
    \DeduceC{$!A(a,c)$}
    \RightLabel{\Intro{\lforall}}
    \UnaryInfC{$\lforall[x][!A(x,c)]$}
    \DisplayProof
    \]
    विगमन गृहीतकानुसार
    $\Subst{\delta'}{t}{c}$
    ही $!A(a,t)$ ची
    नियमित !!{proof} आहे.
    $a$ हे~$\delta$
    चे आयगेन चर असल्याने
    $a$ हे~$t$ मध्ये येत नाही.
    म्हणून
    $\Subst{\delta}{t}{c}$
    मधील अंतिम अनुमान
    \[
    \AxiomC{}
    \RightLabel{$\Subst{\delta'}{t}{c}$}
    \DeduceC{$!A(a,t)$}
    \RightLabel{\Intro{\lforall}}
    \UnaryInfC{$\lforall[x][!A(x,t)]$}
    \DisplayProof
    \]
    योग्य आहे: $t$
    मध्ये~$a$ नसल्याने
    निष्कर्षात किंवा कोणत्याही
    उघड्या गृहीतकात~$a$
    येत नाही.
    \item अंतिम अनुमान
    $\Elim{\lforall}$
    असेल, तर~$\delta$
    चे रूप असे आहे:
    \[
    \AxiomC{}
    \RightLabel{$\delta'$}
    \DeduceC{$\lforall[x][!A(x,c)]$}
    \RightLabel{\Elim{\lforall}}
    \UnaryInfC{$!A(s(c),c)$}
    \DisplayProof
    \]
    विगमन गृहीतकानुसार
    $\Subst{\delta'}{t}{c}$
    ही $\lforall[x][!A(x,t)]$
    ची नियमित !!{proof}
    आहे. (निष्कर्षातील
    पद~$s(c)$ मध्ये~$c$
    असूही शकतो किंवा नसूही
    शकतो.) आता
    $\Subst{\delta}{t}{c}$
    मधील अंतिम अनुमान
    \[
    \AxiomC{}
    \RightLabel{$\Subst{\delta'}{t}{c}$}
    \DeduceC{$\lforall[x][!A(x,t)]$}
    \RightLabel{\Elim{\lforall}}
    \UnaryInfC{$!A(s(t),t)$}
    \DisplayProof
    \]
    योग्य आहे, आणि
    $!A(s(t),t) \ident
    \Subst{!A(s(c),c)}{t}{c}$.
    \Intro{\lexists} चा
    प्रसंगही असाच आहे.
  \item अंतिम अनुमान
    $\Elim{\lexists}$
    असेल, तर~$\delta$
    चे रूप असे आहे:
    \[
    \AxiomC{}
    \RightLabel{$\delta_1'$}
    \DeduceC{$\lexists[x][!A(x,c)]$}
    \AxiomC{$\Discharge{!A(a,c)}{x}$}
    \RightLabel{$\delta_2'$}
    \DeduceC{$!C$}
    \DischargeRule{\Elim{\lexists}}{x}
    \BinaryInfC{$!C$}
    \DisplayProof
    \]
    विगमन गृहीतकानुसार
    $\Subst{\delta_1'}{t}{c}$
    आणि
    $\Subst{\delta_2'}{t}{c}$
    या अनुक्रमे
    $\lexists[x][!A(x,t)]$
    ची आणि उघड्या
    गृहीतक~$!A(a,t)$ पासून
    $\Subst{!C}{t}{c}$
    ची नियमित !!{proof}s
    आहेत. $a$ हे~$\delta$
    चे आयगेन चर असल्याने
    $a$ हे~$t$ मध्ये येत नाही.
    म्हणून
    $\Subst{\delta}{t}{c}$
    मधील अंतिम अनुमान
    \[
    \AxiomC{}
    \RightLabel{$\Subst{\delta_1'}{t}{c}$}
    \DeduceC{$\lexists[x][!A(x,t)]$}
    \AxiomC{$\Discharge{!A(a,t)}{x}$}
    \RightLabel{$\Subst{\delta_2'}{t}{c}$}
    \DeduceC{$\Subst{!C}{t}{c}$}
    \DischargeRule{\Elim{\lexists}}{x}
    \BinaryInfC{$\Subst{!C}{t}{c}$}
    \DisplayProof
    \]
    योग्य आहे:
    $\Subst{!A(x,t)}{a}{x}
    \ident \Subst{!A(a,c)}{t}{c}$;
    निष्कर्ष~$\Subst{!C}{t}{c}$
    किंवा कोणत्याही उघड्या
    गृहीतकात~$a$ येत नाही.
  \end{enumerate}
\end{proof}

\begin{prop}\ollabel{prop:regular}
$\delta$ ही \Log{N1c}
किंवा \Log{N1i} मधील
!!a{proof} असेल,
तर त्याच रचनेची
(विशेषतः त्याच
!!{height} ची)
नियमित !!{proof}~$\delta'$
असते; ती~$\delta$
पेक्षा केवळ आयगेन
चरांच्या बदलामुळे
भिन्न असते.
\end{prop}

\begin{proof}
या सिद्धतेपुरते,
~$\delta$ मधील एखादे
आयगेन-चर अनुमान
\emph{स्वच्छ} म्हणू,
जर त्याचे आयगेन चर
इतर कोणत्याही अनुमानाचे
आयगेन चर नसेल आणि
त्या अनुमानाच्या निकटच्या
उप-!!{proof} व्यतिरिक्त
~$\delta$ मध्ये
कुठेही येत नसेल;
अन्यथा ते \emph{अस्वच्छ}
म्हणू. ~$n$ ही~$\delta$
मधील अस्वच्छ आयगेन-चर
अनुमानांची संख्या असू द्या.
~$n$ वर विगमन करून
~$\delta$ चे अपेक्षित
नियमित !!{proof}~$\delta'$
मध्ये रूपांतर करता येते,
हे दाखवू. $n=0$ असेल,
तर~$\delta$ मधील सर्व
आयगेन-चर अनुमाने स्वच्छ
आहेत; म्हणून~$\delta$
नियमित आहे आणि आणखी
काही सिद्ध करायचे नाही.
अन्यथा सर्वात वरचे
अस्वच्छ आयगेन-चर अनुमान
निवडा; उदाहरणार्थ
\Intro{\lforall}.
त्यात संपणारी~$\delta$
ची उप-!!{proof}~$\delta_1$
या रूपाची आहे:
    \[
    \AxiomC{}
    \RightLabel{$\delta_2$}
    \DeduceC{$!A(c)$}
    \RightLabel{\Intro{\lforall}}
    \UnaryInfC{$\lforall[x][!A(x)]$}
    \DisplayProof
    \]
~$c$ हे आयगेन चर असल्याने
ते~$\lforall[x][!A(x)]$
मध्ये किंवा उघड्या गृहीतकात
येत नाही. ~$\delta_2$
नियमित आहे: निवडलेल्या
अनुमानाच्या वर असलेली
सर्व आयगेन-चर अनुमाने
स्वच्छ आहेत. ~$a$ हा
~$\delta$ मध्ये कुठेही
न येणारा !!a{constant}
असू द्या.
\olref{lem:subst-regular}
नुसार
$\Subst{\delta_2}{a}{c}$
ही $!A(a)$ ची नियमित
सिद्धता आहे.
आयगेन चराच्या अटीने
~$\delta_2$ च्या कोणत्याही
उघड्या गृहीतकात~$c$
नसतो; म्हणून
~$\Subst{\delta_2}{a}{c}$
च्या उघड्या गृहीतकात
~$a$ नसतो. त्यामुळे
~$\Intro{\lforall}$
लावता येतो.
~$\delta$ मधील
$\delta_1$ च्या जागी
पुढील सिद्धता~$\delta_1'$
ठेवा:
    \[
    \AxiomC{}
    \RightLabel{$\Subst{\delta_2}{a}{c}$}
    \DeduceC{$!A(a)$}
    \RightLabel{\Intro{\lforall}}
    \UnaryInfC{$\lforall[x][!A(x)]$}
    \DisplayProof
    \]
यात एक अस्वच्छ आयगेन-चर
अनुमान स्वच्छ केले आहे,
आणि फक्त आयगेन चर~$c$
च्या जागी~$a$ आले आहे.
परिणामी सिद्धतेत
$n-1$ अस्वच्छ आयगेन-चर
अनुमाने राहतात; म्हणून
विगमन गृहीतकाने अपेक्षित
निष्कर्ष मिळतो.
\end{proof}

\begin{prob}
\olref{prop:regular} च्या
सिद्धतेत सर्वात वरचे
आयगेन-चर अनुमान
~\Elim{\lexists} असते
तो प्रसंग वर्णन करा.
\end{prob}

नुकतेच सिद्ध केलेले
विधान आपण पुढे स्पष्टपणे
वापरणार नाही; पण चर्चेतील
सर्व !!{proof}s नियमित
आहेत असे मानू.
तसे नसल्यास आयगेन
चरांची नावे बदलून त्या
नियमित करता येतात.

\Olref{lem:subst-regular}
आणि \olref{prop:regular}
ही विधाने \Log{N2c}
आणि~\Log{N2i}
यांनाही लागू होतात.
त्यांच्या सिद्धता
सरावासाठी ठेवतो.

\end{document}
